% Based on templates/lecture-example.tex; actual lecture SC3000-L09.
% Source: Week 10 - First-order Logic.pptx, original slide order.
% Lecture examples are paraphrased; derivations and the minimal Prolog version
% are explanatory reconstructions, not a transcription of the slide 56 program.
\begingroup
\renewcommand{\recordid}[1]{}
\renewcommand{\relatedtopic}[2]{\hyperref[topic:#1]{#2}}

\exercisegroup{SC3000-L09-EX}{第 09 讲：First-Order Logic}

\textbf{日期与来源：}2026-10-19；Week 10 -- First-order Logic，pp.~10--12、31、34--35、49--56。题意按课件例子整理，补全推导；代码为亲属例子的最小 Prolog 改写。

\exercisequestion{SC3000-L09-E01}{Socrates 三段论：一般规则如何用于具体对象}

\textbf{对应知识点：}\relatedtopic{SC3000-L09-T01}{结构化表示}；\relatedtopic{SC3000-L09-T05}{量词推理}。

\textbf{题目。}课件 pp.~10--12 给出“Socrates 是人”“所有人都会死”，结论为“Socrates 会死”。说明将三句分别记为独立命题符号的局限，并用 FOL 推导结论。

\begin{workedsolution}
若三个句子分别编码为 $P,Q,R$，则 $P=Q=\mathrm T,R=\mathrm F$ 仍是命题逻辑允许的赋值。因此仅有 $P,Q$ 不蕴含 $R$；还需要额外写入连接它们的规则。

FOL 将性质和对象分开表示：
\[
 Man(Socrates),\qquad
 \forall x\,[Man(x)\Rightarrow Mortal(x)].
\]
先以 $x\leftarrow Socrates$ 作全称实例化，再作 MP：
\[
 Man(Socrates)\Rightarrow Mortal(Socrates),\quad
 Man(Socrates)\ \vdash\ \boxed{Mortal(Socrates)}.
\]
同一条全称规则也能用于 $Plato$；推理依靠公式中共享的谓词与对象结构，不依靠系统理解自然语言中的人名。
\end{workedsolution}

\exercisequestion{SC3000-L09-E02}{狗主人与共同敌人：从嵌套条件读出句意}

\textbf{对应知识点：}\relatedtopic{SC3000-L09-T02}{量词与作用域}；\relatedtopic{SC3000-L09-T03}{等价改写}。

\textbf{题目。}按 pp.~31、34--35，改写并解释：
\begin{align*}
 &(1)\quad Owns(x,y)\Rightarrow[Dog(y)\Rightarrow AnimalLover(x)],\\
 &(2)\quad [Dog(x)\land Enemy(z,x)]\Rightarrow[Dog(y)\Rightarrow Enemy(z,y)].
\end{align*}
这些公式作为一般规则时，所有出现的变量均取全称量化；约定 $Enemy(z,x)$ 为“$z$ 是 $x$ 的敌人”。

\begin{workedsolution}
首先验证共同变换：
\begin{align*}
 A\Rightarrow(B\Rightarrow C)
 &\equiv\neg A\lor(\neg B\lor C)\\
 &\equiv\neg(A\land B)\lor C
 \equiv(A\land B)\Rightarrow C.
\end{align*}
因此第一式为
\[
 \forall x\forall y\,
 [Owns(x,y)\land Dog(y)\Rightarrow AnimalLover(x)].
\]
它说“拥有狗的人是爱动物的人”，没有反向断言每个爱动物的人都养狗，也没有断言狗主人一定存在。

第二式为
\[
 \forall x\forall y\forall z\,
 [Dog(x)\land Dog(y)\land Enemy(z,x)\Rightarrow Enemy(z,y)].
\]
任取两只狗，第一只的任何敌人也是第二只的敌人。由于 $x,y$ 任意，交换二者后也成立，故两只狗的敌人集合相同。这不保证敌人存在；所有狗的敌人集合都为空时，规则同样成立。
\end{workedsolution}

\exercisequestion{SC3000-L09-E03}{亲属关系：定义、变量绑定与 Prolog 查询}

\textbf{对应知识点：}\relatedtopic{SC3000-L09-T04}{亲属知识库}；\relatedtopic{SC3000-L09-T05}{逻辑程序}。

\textbf{题目。}课件 pp.~49--51 的论域为人，$Parent(x,y)$ 表示 $x$ 是 $y$ 的家长。已知 $Female(Elizabeth)$、$Parent(Elizabeth,Charles)$ 和 $Parent(Charles,William)$。用一般规则定义父母函数、子女、祖父母和兄弟姐妹关系，并回答 $Grandparent(Elizabeth,William)$ 与 $Child(William,x)$ 两个查询。

\begin{workedsolution}
以下规则均对自由出现的 $x,y$ 全称量化。母亲和父亲函数采用课件的简化亲属模型，每个输入对应一个确定对象：
\begin{align*}
 FatherOf(x)=y&\Leftrightarrow Parent(y,x)\land Male(y),\\
 MotherOf(x)=y&\Leftrightarrow Parent(y,x)\land Female(y),\\
 Parent(x,y)&\Leftrightarrow Child(y,x),\\
 Grandparent(x,y)&\Leftrightarrow
   \exists z\,[Parent(x,z)\land Parent(z,y)],\\
 Sibling(x,y)&\Leftrightarrow
   [x\ne y\land\exists z\,(Parent(z,x)\land Parent(z,y))].
\end{align*}
Sibling 的定义要求不同的人共享至少一个家长；$x\ne y$ 防止把自己列为自己的兄弟姐妹。

\begin{center}
% Original diagram based on Week 10 FOL slides 50--51.
\begin{tikzpicture}[
  >=Latex,
  person/.style={draw=noteBlue,rounded corners=2pt,
    minimum height=0.65cm,minimum width=2.1cm,font=\small},
  relation/.style={->,thick,noteBlue}
]
  \node[person] (e) at (0,0) {Elizabeth};
  \node[person] (c) at (4.5,0) {Charles};
  \node[person] (w) at (9,0) {William};
  \draw[relation] (e) -- node[above,font=\small]{Parent} (c);
  \draw[relation] (c) -- node[above,font=\small]{Parent} (w);
  \draw[relation,dashed] (e.south) -- ++(0,-0.7)
    -- node[below,font=\small]{Grandparent} ([yshift=-0.7cm]w.south)
    -- (w.south);
\end{tikzpicture}


\small 依据课件 pp.~50--51 的事实重绘；实线是已知关系，虚线是推导结果。
\end{center}

由两条 Parent 事实作合取引入，得到
\[
 Parent(Elizabeth,Charles)\land Parent(Charles,William).
\]
以 Charles 为见证作存在引入，再应用祖父母规则的一侧，得到
\[
 \exists z\,[Parent(Elizabeth,z)\land Parent(z,William)]
 \ \vdash\ \boxed{Grandparent(Elizabeth,William)}.
\]
因此第一个查询可回答 yes。同理由子女规则可确认 William 是 Charles 的孩子，第二个查询返回绑定 $\{x/Charles\}$。查询方向不能倒置；已知事实也不保证 Charles 是唯一可能满足查询的家长。

\textbf{Prolog 最小改写。}下列程序取祖父母定义的充分条件方向，演示课件 pp.~55--56 所述的事实、规则和查询方式：
\begin{verbatim}
parent(elizabeth, charles).
parent(charles, william).
grandparent(X, Y) :- parent(X, Z), parent(Z, Y).
child(Y, X) :- parent(X, Y).

?- grandparent(elizabeth, william).
true.
?- child(william, X).
X = charles.
\end{verbatim}
第一条查询通过 $Z=charles$ 连通两代关系；第二条查询返回满足条件的绑定。
\end{workedsolution}
\endgroup
